\subsection{Negative odd roots at every nonunit inner length}\label{proof:Q5-C080}
\paragraph{Q80.0 Original objective and exact new claims}

The ambient problem includes the P5, T5, and K1, 4 skeletons of Q5.
The following estimates treat the specified root roles.

Fix the finite probability space and kernel throughout. Write

\fitdisplay{
Jv=(Ev)1,\quad A=J+K,\quad T=J-K=A^*,\quad
 g=K1,\quad p=1+g,\quad q=1-g,
}{}
\fitdisplay{
f_n=A^{n}\one,\quad \bar f_n=T^{n}\one,\quad
 h_n=T f_{n-1},\quad W_n=A \bar f_{n-1},\quad f_0=\bar f_0=1.
}{}

Take the strictly ordered original odd residual root table

\fitdisplay{
a<u_1<v_1<\cdots<u_k<v_k,\qquad k\ge1,
 \qquad P=\prod_{i=1}^k W_{u_i}\bar f_{v_i}.
}{Q80.0.1}

The terminal has original residual arms (1, b), with every integer \(b\ge 2.\)
For every nonunit inner length \(L\ge 2,\allowbreak \) the actual ordinary input is

\fitdisplay{
X_{b,L}=T^L(qf_b),\qquad Y_{a,b,L}=W_aP X_{b,L}.
}{Q80.0.2}

Theorem T085.1. If \(a\ge 3,\allowbreak \) all original hosts and all the independent
finite lengths in (Q80.0.1)--(Q80.0.2) satisfy

\fitdisplay{
\boxed{EY_{a,b,L}^{8}\le e^{-97d/12},\qquad
 I_q(T^M Y_{a,b,L})^3\le e^{-49d/228}\quad(M\ge1),}
}{Q80.0.3}

where \(d=\lVert K1\rVert _2^2\) and \(I_q(U)=\E[q^{-1/3}|U|^{8/3}]\).
For \(a=3\) the ordinary exponent improves to \(403/48,\allowbreak \) and the all-M
weighted exponent improves to \(481/912.\)

Theorem T085.2. If \(a=2\) and \(u_1\ge 4,\allowbreak \) every \(b\ge 2,\allowbreak  L\ge 2,\allowbreak  M\ge 1\) satisfies

\fitdisplay{
\boxed{EY_{2,b,L}^{8}\le e^{-479d/60},\qquad
 I_q(T^M Y_{2,b,L})^3\le e^{-131d/1140}.}
}{Q80.0.4}

The factors in (Q80.0.2)
are the complete actual root input, and \(T^M\) is its original outer path.


Theorem T085.3, exact all-length union. For every odd root table
(Q80.0.1) with \(a\ge 3,\allowbreak \) terminal residual (1, b), \(b\ge 2,\allowbreak \) and all independent
positive original L, M, choose the following directions from L alone:


\begingroup\small
\setlength{\tabcolsep}{3pt}
\setlength{\LTleft}{0pt}\setlength{\LTright}{\fill}
\begin{longtable}{@{}p{0.1200\linewidth}p{0.8000\linewidth}@{}}
\caption{Negative odd roots at every nonunit inner length: parameter values.}\label{tab:parameters:q5-c080:1}\\
\toprule
Entry & Complete ordinary root input before the outer path \\
\midrule\endfirsthead
\toprule
Entry & Complete ordinary root input before the outer path \\
\midrule\endhead
\bottomrule\endfoot
\(L=1\) & \(Y=f_aP T(p\bar f_b)\). \\
\addlinespace[2pt]
\(L\ge 2\) & \(Y=W_aP T^L(qf_b)\). \\
\end{longtable}
\endgroup


Then, including equal-arm factors at their roots by positivity,

\fitdisplay{
\boxed{I_q(T^M Y)^3\le e^{-25d/912}\le1.}
}{Q80.0.5}

All claimed root and terminal directions are literal original directions:

\begin{description}[style=nextline,leftmargin=1.6em]
\item[{\(Root a-arm\)}] Nonunit row of (Q80.0.2): First edge away from root, all later edges toward its first neighbor;\par\noindent factor \(W_a\).
\item[{\(Root u_i-arm\)}] Nonunit row of (Q80.0.2): Same actual reversed bent pattern;\par\noindent factor \(W_{u_i}\).
\item[{\(Root v_i-arm\)}] Nonunit row of (Q80.0.2): \(Entirely toward root; factor \bar f_{v_i}\).
\item[{Terminal unit arm}] Nonunit row of (Q80.0.2): Toward terminal;\par\noindent factor q.
\item[{Terminal b-arm}] Nonunit row of (Q80.0.2): Entirely away from terminal;\par\noindent factor \(f_b\).
\item[{Inner connector}] Nonunit row of (Q80.0.2): All original L edges from terminal toward root.
\item[{Outer connector}] Nonunit row of (Q80.0.2): All original M edges from root toward receiving root.
\end{description}

For \(L=1,\allowbreak \) replace the a-arm by its entirely outward orientation \(f_a,\allowbreak \)
and the terminal by the actual unit-outward, b-arm-inward pattern \(p\bar f_b.\)
The extra pairs P and both core-path directions remain as displayed.
The original integer L, not any host parameter, determines this choice.



\paragraph{Q80.1 Original calculus and one finite mask for all extra pairs}

Delete zero-mass atoms. Set

\fitdisplay{
r=Kg,\quad s_0=Kr,\quad t=Ks_0,\quad
 \delta_j=\|K^{j}\one\|_2^2,\quad d=\delta_1,\quad
 \theta=13/32,\quad\kappa=19/32,\quad\tau=d-\delta_2.
}{}

The C014 calculus gives

\fitdisplay{
0\le d\le1/3,\quad d^2\le \delta_2\le\theta d,\quad
 \|K\|^2\le\theta,\quad \delta_{j+1}\le\theta \delta_j,
 \quad |Kv(x)|\le\|v\|_2.
}{Q80.1.1}

The original two-row estimate is \(|K^2v(x)|^2\le \lVert v\rVert _2^2/3.\)
It is not an operator-norm bound \(\lVert K\rVert ^2\le 1/3.\)
The path convolution uses the finite quantities

\fitdisplay{
\beta_1=\beta_2=0,\qquad
 \beta_n=\sum_{j=1}^{\lfloor(n-1)/2\rfloor}(-1)^{j-1}\delta_j,
 \quad\tau\le\beta_n\le d\quad(n\ge3),
}{}
\fitdisplay{
\sum_{j=0}^n K^{j}\one=f_n+\sum_{j=3}^n\beta_jf_{n-j}.
}{Q80.1.2}

The same identity holds after complete reversal \(K->-K.\) Every path
message here is nonnegative. In fact p, \(q\ge \mu(x)>0\)
on each retained atom. Actual equal pairs obey \(f_l h_l\le 1\) and,
after complete reversal, \(W_l \bar f_l\le 1.\)

At \(d=0,\allowbreak  g=0\) while K may be nonzero. Nevertheless \(A\one=T\one=\one,\allowbreak \)
and successive leaf integrations make every tree
message here equal to one. Every displayed final bound
holds directly. All later divisions by d are restricted to \(d>0.\)

\paragraph{Q80.1.1 The exact actual pair loss and original finite mask}

C067 proves, including the separate finite \(u=4\) case,

\fitdisplay{
h_u f_v\le\exp\left(\sum_{j=u}^{v-1}K^jg-3d\right)
 \qquad(4\le u<v).
}{Q80.1.3}

We use this only with \(u\ge 4.\) Full reversal on the same original
probability space gives

\fitdisplay{
P\le e^{S-3kd},\qquad
 S=-\sum_{i=1}^k\sum_{j=u_i}^{v_i-1}(-K)^jg=KV_S,
}{}
\fitdisplay{
V_S=U(-K)g,\qquad
 U(z)=\sum_{i=1}^k\sum_{j=u_i}^{v_i-1}z^{j-1}.
}{Q80.1.4}

The intervals are disjoint. Thus U is a finite polynomial whose
coefficients are zero or one, independently of the number of arms,
the gaps, and the largest original length. 

For every finite [0, 1]-coefficient polynomial Q and \(z=i \sqrt{x}\),
\(0\le x\le \theta,\allowbreak \) separate the four sign classes. Their finite sums give

\fitdisplay{
|\Re Q(z)|\le(1-x^2)^{-1},\qquad
 |\Im Q(z)|\le\sqrt x(1-x^2)^{-1},
}{}
\fitdisplay{
\|Q(K)\|^2\le\frac{1+\theta}{(1-\theta^2)^2}
 =\frac{32768}{16245}<\frac{100}{49}.
}{Q80.1.5}

The comparison is \(32768*49=1605632<1624500=16245*100.\)
Only a finite polynomial is evaluated; the geometric sums are upper
bounds for finite positive subcollections.

If \(u_1\ge 4,\allowbreak \) U begins in degree at least three. If \(u_1\ge 5\) it begins
in degree at least four. Consequently

\fitdisplay{
\|V_S\|\le
 \begin{cases}
 (3/8)\sqrt d,&u_1\ge4,\\
 (1/4)\sqrt d,&u_1\ge5.
 \end{cases}
}{Q80.1.6}

For the first, (\(10/7)\theta^{3/2}<585/1568<3/8,\allowbreak \) using
\(\sqrt{\theta}<9/14.\) For the second,
(\(10/7)\theta^2=845/3584<1/4.\) Original row Cauchy applies only
to the outermost K in \(S=KV_S,\allowbreak \) not to a compressed or spectral kernel.

The positive contributions to ES have \(j=3,\allowbreak \) 7, 11,... and values
\(\delta_2,\allowbreak  \delta_4,\allowbreak  \delta_6,\allowbreak ...\) . Since \(u_1\ge 4\) excludes the first such term,

\fitdisplay{
ES\le\frac{\delta_4}{1-\theta^2}
 \le\frac{2197}{27360}d<\frac d{12}.
}{Q80.1.7}

Equations (Q80.1.4)--(Q80.1.7) hold for every finite k and every original
set of gaps. Unequal pairs are paid by these inequalities, rather
than deleted from the tree or assumed pointwise at most one.

\paragraph{Q80.2 A precise envelope extension with the same original fluctuation}\label{part:Q80.2}

Let \(c_*=(1+\theta)^2=2025/1024<2.\) The exact C071/C075 exponential
tool says that if \(Z=KV,\allowbreak  \lVert V\rVert ^2\le D,\allowbreak \) and \(M\ge max(EZ,\allowbreak \) 0), then

\fitdisplay{
\log Ee^{8Z/B}\le\frac{8M}{B}+\frac{8c_*D}{B^2}
 \qquad(B>0).
}{Q80.2.1}

Its hypotheses are the original range \(|Z|^2\le D\) and original raw
second moment \(EZ^2\le \theta\) D. It covers both signs of EZ.
For \(Z=K^2V_2\) with \(\lVert V_2\rVert ^2\le D_2,\allowbreak \) the original two-row version is

\fitdisplay{
\log Ee^{8Z/B}\le\frac{8M}{B}+\frac{6D_2}{B^2}.
}{Q80.2.2}

Writing the same two-row variable as \(K(KV_2)\) gives a one-row
primitive with squared norm at most theta \(D_2.\) Its scalar cost
increases by at most

\fitdisplay{
(8c_*\theta-6)D_2/B^2
 =\frac{1749}{4096}D_2/B^2<\frac7{16}D_2/B^2.
}{Q80.2.3}

Apply these scalar lemmas with their stated hypotheses.

Lemma T085.4. Suppose an actual nonnegative base \(Y_0\) has specified
\(C_0,\allowbreak  B>0,\allowbreak \) a real loss ell, a vector V, \(M\ge max(EKV,\allowbreak \) 0), and \(D\ge \lVert V\rVert ^2,\allowbreak \)
with

\fitdisplay{
Y_0\le C_0\left(1+\frac{KV-\ell}{2B}\right)^2,
 \qquad 1+\frac{KV-\ell}{2B}\ge0,
 \quad D\le\eta^2B^2d.
}{Q80.2.4}

Let its complete specified one-row scalar expression be

\fitdisplay{
\mathcal E_0=8\log C_0-8\ell/B+8M/B+8c_*D/B^2.
}{Q80.2.5}

For the actual P in (Q80.1.4), if \(\lVert V_S\rVert \le \epsilon \sqrt{d}\), then

\fitdisplay{
\boxed{E(PY_0)^8\le
 \exp\left[\mathcal E_0+
 \left\{16(2\eta\epsilon+\epsilon^2)+\frac23-24k\right\}d
 \right].}
}{Q80.2.6}

Proof. Positivity of the square base and \(\log(1+x)\le x\) give

\fitdisplay{
(PY_0)^8\le C_0^8
 \exp\{8(KV-\ell)/B+8S-24kd\}.
}{}

At a zero square base this follows by continuity. Combine the two
remainders before taking expectations:

\fitdisplay{
KV+BS=K(V+BV_S).
}{}

Its mean is bounded by the nonnegative M'=M+Bd/12. Its specified
primitive norm can be bounded by

\fitdisplay{
D'=D+2\sqrt D\,B\epsilon\sqrt d+B^2\epsilon^2d
 \le D+B^2(2\eta\epsilon+\epsilon^2)d.
}{}

Use (Q80.2.1) on this same original function, and retain (Q80.2.5) with
its own originally paid D. The extra mean cost is 2d/3; the extra
norm cost is at most \(8c_*(2\eta \epsilon+\epsilon^2)d,\allowbreak \) hence at most
16(2eta \(\epsilon+\epsilon^2)d.\) This proves (Q80.2.6).
 In particular, a
looser norm normalization is used only for the increment; it does
not replace D inside an already paid scalar expression. QED.

\paragraph{Q80.3 The complete actual base table and its exact normalizations}

We use the actual shortest-two construction of C071,
with C075's additional all-domain \(b=2,\allowbreak  L=2\) payment and its shortest-three
comparison. C076/T082 records the following full normalization.
We reproduce all rows needed for the new estimates.

Since \(W_2=Aq=p-r,\allowbreak \) the exact positive terminal masks give

\fitdisplay{
0<X_{b,L}\le C[Bq+Br+Z-\ell],\qquad
 W_2X_{b,L}\le CB\left(1+\frac{Z-\ell}{2B}\right)^2.
}{Q80.3.1}

The bracket is half the sum of two actual nonnegative upper factors.
Thus the sign condition in (Q80.2.4) is satisfied, even before subsequent
root comparisons.

For \(b=2\) put \(\gamma=d-\delta_2/2\) and \(B_2=1-\gamma.\) For \(b\ge 3\) put

\fitdisplay{
R=KG_{b-1}(K)g,\quad G_n(z)=\sum_{j=0}^{n-1}z^j,
 \quad m_b=1-\beta_b/2,\quad k_b=1/m_b,
}{}
\fitdisplay{
t_0=\beta_{b+1},\quad e=d-\beta_{b+2},\quad V_R=ER^2,
 \quad\rho=E(rR),\quad
 \gamma=k_bt_0-k_b^2V_R/2,\quad a_1=k_be,
 \quad B_*=\gamma+a_1.
}{Q80.3.2}

The subscript on \(k_b\) distinguishes it from the number k of root pairs.
The exact terminal toolkit gives \(1\le k_b\le 6/5,\allowbreak  ER=-t_0,\allowbreak  E(gR)=-e\le 0,\allowbreak  V_R=\delta_2+\delta_3\) for \(b=3\) and \(V_R\le \delta_2\) for \(b\ge 4,\allowbreak  0\le a_1\le 39d/80,\allowbreak  \gamma\le 6d/5,\allowbreak \) and \(B_*\le 3d/2.\)

For the finite polynomials in the table set \(z=-K\) and

\fitdisplay{
P_n(z)=
 \begin{cases}
 G_n(z)+z^{n+1},&b=2,\\
 G_n(z)+k_bz^{n+1}G_{b-1}(-z),&b\ge3.
 \end{cases}
}{Q80.3.3}

The entire finite table, including every short connector, is:


\begingroup\small
\setlength{\tabcolsep}{3pt}
\setlength{\LTleft}{0pt}\setlength{\LTright}{\fill}
\begin{longtable}{@{}>{\raggedright\arraybackslash}p{0.1200\linewidth}>{\raggedright\arraybackslash}p{0.2667\linewidth}>{\raggedright\arraybackslash}p{0.2667\linewidth}>{\raggedright\arraybackslash}p{0.2667\linewidth}@{}}
\caption{Finite masks and row-form scalar bounds for the terminal parameters \((b,L)\).}\label{tab:parameters:q5-c080:3}\\
\toprule
\((b,L)\) & C, B and ell & Same original fluctuation Z & Original row-form scalar upper \\
\midrule\endfirsthead
\toprule
\((b,L)\) & C, B and ell & Same original fluctuation Z & Original row-form scalar upper \\
\midrule\endhead
\bottomrule\endfoot
2, 2 & \(C=1; B=B_2; \ell=0\); & \(K^2(\gamma 1+r)\); & \(-d/8\). \\
\addlinespace[2pt]
2, 3 & \(C=1; B=B_2; \ell=\tau\); & \(-K^2(g+s_0)\); & \(-3d/4\). \\
\addlinespace[2pt]
2, 4 & \(C=1; B=1-\gamma-\tau; \ell=d(1-\gamma)\); & \(K^2[\tau 1-(1-\gamma)g+r+t]\); & \(-d\). \\
\addlinespace[2pt]
2, 5 & \(C=1; B=1-2\tau; \ell=\tau\); & \(K[\tau+zP_3(z)]g\); & \(-2d/3\). \\
\addlinespace[2pt]
\(2,\allowbreak L\ge 6\) & \(C=1; B=1-2\tau; \ell=\tau\); & \(K[\tau+zP_{L-2}(z)]g\); & \(-4d/5\). \\
\addlinespace[2pt]
\(b\ge 3,\allowbreak 2\) & \(C=m_b^2; B=1-\gamma; \ell=a_1\); & \(K^2(\gamma 1+k_bR)\); & \(-d/8\). \\
\addlinespace[2pt]
\(b\ge 3,\allowbreak 3\) & \(C=m_b^2; B=1-B_*; \ell=d-k_b \rho\); & \(K[a_1+zP_1(z)]g\); & \(-d/10\). \\
\addlinespace[2pt]
\(b\ge 3,\allowbreak 4\) & \(C=m_b^2; B=1-41d/80; \ell=\kappa d\); & \(K[41d/80+zP_2(z)]g\); & \(-8d/25\). \\
\addlinespace[2pt]
\(b\ge 3,\allowbreak L\ge 5\) & \(C=B=1-d/2; \ell=d/2\); & \(K[d/2+zP_{L-2}(z)]g\); & \(-2d/5\). \\
\end{longtable}
\endgroup


Here the original row-form scalar is

\fitdisplay{
8\log(CB)-8\ell/B+8M/B+
 \begin{cases}
 6D_2/B^2,&Z=K^2V_2,\\
 8c_*D_1/B^2,&Z=KV_1.
 \end{cases}
}{Q80.3.4}

The \(D_j\) in (Q80.3.4) are the specified bounds for those particular vectors.
Each paid scalar expression retains its own shared terminal parameters.
These bounds give (Q80.3.4) directly.

Nonnegative mean uppers M in the displayed order are
\(\delta_2-\gamma\) d, 0, (theta kappa-tau)d, (theta-tau)d for both \(b=2\) tails,
\(max(k_b\) rho-gamma d, 0), and (theta-delta)d in the last three rows,
where \(\delta=a_1,\allowbreak \) 41d/80, d/2 respectively. They bound the original Z
in each row and satisfy \(M\ge 0\) over the entire original domain.

\paragraph{Q80.3.1 Two-row vectors and conversion costs}

The four two-row vector squares have the following whole-domain
normalizations. These upper bounds are used for increments and norm
comparisons, not substituted into a scalar payment that used a sharper
same-parameter expression:

\fitdisplay{
\begin{array}{c|c}
 (b,L)&D_2/(B^2d)\\\hline
 (2,2)&\theta\,324/169<1\\
 (2,3)&324/169<2\\
 (2,4)&217/81<25/9\\
 (b\ge3,2)&540/289<2.
 \end{array}
}{Q80.3.5}

The original lower bounds are \(B_2\ge 13/18,\allowbreak  1-\gamma-\tau\ge 1/2,\allowbreak \) and \(1-\gamma\ge 17/25\) for \(b\ge 3.\) The vector squares
are bounded respectively by \(\delta_2,\allowbreak  d-2\delta_2+\delta_3\le d,\allowbreak \)
[(\(1-\gamma)^2+4/27\)]d, and \(k_b^2 V_R.\)
For \(b=2,\allowbreak  L=4,\allowbreak \) use \(1-\gamma=B+\tau\) and \(\tau/B\le 4/9,\allowbreak \) giving
(\(1+4/9)^2+16/27=217/81.\) For \(b\ge 3,\allowbreak  L=2,\allowbreak  \theta(1+\theta)=585/1024<3/5\) gives the upper \(540/289.\)

After the exact two-row to one-row conversion (Q80.2.3), the squared
primitive normalization is multiplied by theta, while the specified
scalar upper increases by less than (\(7/16)D_2/B^2.\)

\paragraph{Q80.3.2 Every one-row tail}

For \(b=2,\allowbreak  L\ge 5\) the old squared primitive is bounded by A(tau)d,
where

\fitdisplay{
A(\tau)=A_0-2\theta\sigma\tau+\tau^2,
 \quad(A_0,\sigma)=
 \begin{cases}(\theta,1),&L=5,\\
 (\theta-7\theta^2/15,\kappa),&L\ge6.
 \end{cases}
}{}

Both are bounded above by theta-2theta kappa \(\tau+\tau^2.\)
Its quotient by (\(1-2\tau)^2\) increases for \(0\le \tau\le 2/9:\)
after multiplying the derivative by (\(1-2\tau)^3,\allowbreak \) its numerator is
4theta-2theta kappa+(2-4theta \(\kappa)\tau>0.\)
Its endpoint is \(289/256.\) The other three one-row normalizations are

\fitdisplay{
\frac{D_1}{B^2d}\le
 \begin{cases}
 289/256,&b=2,L\ge5,\\
 481/400,&b\ge3,L=3,\\
 27/32,&b\ge3,L=4,\\
 18/25,&b\ge3,L\ge5.
 \end{cases}
}{Q80.3.6}

The latter three follow from the actual bounds
\(D_1/d\le 351/1280+\delta^2,\allowbreak  \theta(71/63)+\delta^2,\allowbreak  \theta+d^2/4,\allowbreak \)
and the respective \(B\ge 1/2,\allowbreak  4/5,\allowbreak  5/6.\) All four normalizations are
at most \(5/4.\) Thus (Q80.3.5)--(Q80.3.6) include every \(b\ge 2,\allowbreak  L\ge 2,\allowbreak \)
including the finite cases immediately before each tail.

\paragraph{Q80.3.3 The uniform base for \(a\ge 4\)}

C076/T082 Lemma T082.2 proves the following particular one-row
certificate:
for each \(a\ge 2,\allowbreak  b\ge 2,\allowbreak  L\ge 2,\allowbreak \) the actual Y\_0=W\_\(aX_{b,\allowbreak  L}\) admits (Q80.2.4)
with

\fitdisplay{
D\le(9/4)B^2d,\qquad \mathcal E_0\le(9/4)d.
}{Q80.3.7}

Its exact construction for \(a\ge 4\) starts from the rows above. Put

\fitdisplay{
m=a-2,\quad c=1-\tau\ge7/9,\quad
 H=G_m(-K)g-\tau1,\quad h=B/c,
}{}
\fitdisplay{
A_m=45/32\ (m=2),\qquad A_m=1\ (m\ge3).
}{}

The retained actual positive loss tau p gives

\fitdisplay{
W_aX\le CBc\left(1+
 \frac{Z+hK^2H-\ell-h\beta_a}{2B}\right)^2.
}{Q80.3.8}

The bracket is an actual nonnegative arithmetic mean. The exact
Q5-C075 same-kernel comparisons are \(EK^2H\le 0,\allowbreak  \lVert H\rVert ^2\le A_m\) d, and

\fitdisplay{
\|V_2+hH\|^2\le D_2+h^2A_m d,\qquad
 \|V_1+hKH\|^2\le D_1+h^2\theta A_m d.
}{Q80.3.9}

In its original row form, (Q80.3.8) has scalar upper at most -3d/20.
The largest normalized two-row upper is

\fitdisplay{
25/9+3645/1568<36/7,
}{}

with positive difference \(571/14112.\) Conversion adds less than
(\(7/16)(36/7)d=9d/4,\allowbreak \) while its new one-row norm is at most
\(\theta(36/7)=117/56<9/4.\) In an already one-row case its norm is
at most \(5/4+47385/50176<9/4\) and no conversion is needed.
The mean upper remains valid by \(EK^2H\le 0.\) Set \(C_0=CBc,\allowbreak \)
replace ell by ell+h \(\beta_a,\allowbreak \) and use these actual vectors.
This explains exactly the uniform certificate (Q80.3.7) used below.
The finite covariance proof and all actual coefficient end cases
are the specified C075 parts Q80.4 and 6 / C076 \mbox{\hyperref[part:Q80.2]{part Q80.2}} inputs.

\paragraph{Q80.4 A sharper certificate for the actual shortest-three base}

The generic normalization (Q80.3.7) loses too much for the first extra
short arm four. We therefore retain the precise shortest-three identity

\fitdisplay{
W_3=W_2+s_0-d.
}{}

In every original row (Q80.3.1), replace Z by Z+B \(s_0\) and ell by ell+Bd:

\fitdisplay{
W_3X\le CB\left(1+
 \frac{Z+B s_0-\ell-Bd}{2B}\right)^2.
}{Q80.4.1}

The base remains nonnegative, being the arithmetic mean of the actual
\(W_3\) and the positive terminal upper factor. Since \(Es_0=0,\allowbreak \) the old
nonnegative mean upper M still applies. The deterministic improvement
in the scalar expression is exactly -8d.

Lemma T085.5. For every \(b\ge 2,\allowbreak  L\ge 2,\allowbreak \) (Q80.4.1) has a specified
one-row certificate satisfying

\fitdisplay{
\boxed{D\le(81/64)B^2d,\qquad \mathcal E_0\le-13d/16.}
}{Q80.4.2}

Proof, both \(L=2\) terminal rows. The new two-row vector is
\(V_2+B\) g. When \(b=2,\allowbreak \) g is orthogonal to 1 and r. When \(b\ge 3,\allowbreak \)
the additional cross term uses \(E(gR)=-e\le 0.\) Thus one may take
the new vector upper \(D_2+B^2d.\) The old ratio in (Q80.3.5) is below
two, so the new ratio is below three. Its original row-form
scalar is at most

\fitdisplay{
-d/8+6d-8d=-17d/8.
}{}

Conversion to one row costs less than 21d/16, leaving -13d/16.
The new one-row squared primitive is at most 3theta \(B^2d,\allowbreak \)
and \(3\theta=39/32=78/64<81/64.\)

Proof, \(b=2\) and \(L=3\) or 4. For \(L=3\) the vector changes from
-(\(g+s_0)\) to -(gamma \(g+s_0).\) The squared-norm bound follows from the pointwise inequality

\fitdisplay{
(1-x)^2-(\gamma-x)^2=(1-\gamma)(1+\gamma-2x)\ge0.
}{}

For \(L=4\) the new vector is tau 1-tau g+r+t. Orthogonality of g
to 1, r, t and \(\tau<1-\gamma\) again makes its norm no larger.
The normalized old square is at most \(25/9<3.\) The original
scalar is at most -8d, before conversion. After conversion it
is at most -\(8d+21d/16=-107d/16,\allowbreak \) stronger than required.
The squared one-row norm is at most \(\theta(25/9)<81/64.\)

Proof, all remaining one-row cases. The old primitive is
[\(\delta+zP_n(z)\)]g and the new one is [\(\delta+z(P_n(z)-B)\)]g,
with \(z=-K.\) At frequency \(z=i \sqrt{x}\), the difference of squared
moduli is exactly

\fitdisplay{
Bx[B-2\Re P_n(z)].
}{}

The complete C075 finite coefficient lemma gives Re \(P_n\ge 1/2\)
for every relevant finite n, including \(n=1,\allowbreak \) 2 where needed.
Since \(0<B\le 1,\allowbreak \) the difference is nonpositive. The norm therefore
retains (Q80.3.6), at most \(5/4=80/64<81/64.\) Its original scalar
retains the full extra -8d and is at most -8d.
These cases exhaust all b, L and prove (Q80.4.2). QED.

\paragraph{Q80.5 Minimum at least three, every nonunit inner connector}

Take the actual input (Q80.0.2), \(k\ge 1.\)

If \(a\ge 4,\allowbreak \) the first extra short arm satisfies \(u_1\ge a+1\ge 5.\)
Use the precise certificate (Q80.3.7) in Lemma T085.4 with
\(\eta=3/2,\allowbreak  \epsilon=1/4.\) The full exponent is at most

\fitdisplay{
\left[\frac94+
 16\left(2\cdot\frac32\cdot\frac14+\frac1{16}\right)
 +\frac23-24k\right]d
 \le\left(\frac94+13+\frac23-24\right)d
 =-\frac{97}{12}d.
}{Q80.5.1}

If \(a=3,\allowbreak \) the first extra short arm satisfies \(u_1\ge 4.\)
Use the sharper (Q80.4.2), \(\eta=9/8,\allowbreak  \epsilon=3/8.\) The full exponent is

\fitdisplay{
\left[-\frac{13}{16}+
 16\left(2\cdot\frac98\cdot\frac38+\frac9{64}\right)
 +\frac23-24k\right]d
 \le\left(-\frac{13}{16}+\frac{63}{4}+\frac23-24\right)d
 =-\frac{403}{48}d.
}{Q80.5.2}

Since \(403/48>97/12,\allowbreak \) these prove the ordinary part of Theorem T085.1
over the full original host domain. Every additional pair beyond the
first contributes another -24d; its number never increases the finite
mask bound. All original lengths and gaps remain arbitrary.

\paragraph{Q80.6 Minimum two with first extra short arm at least four}

Now take \(a=2,\allowbreak  u_1\ge 4.\) The extra mask has \(\epsilon=3/8.\)
Each actual shortest-two row of (Q80.3.4)--(Q80.3.6) retains its own paid scalar cost.
For a normalized primitive upper eta, the increment in (Q80.2.6) is

\fitdisplay{
16(2\eta\cdot3/8+9/64)d+(2/3)d-24kd.
}{Q80.6.1}

The following six rows exhaust every \(b\ge 2,\allowbreak  L\ge 2.\) The column \(E_0/d\)
is an upper for the complete specified one-row scalar after any
necessary two-row conversion.


\begingroup\small
\setlength{\tabcolsep}{3pt}
\setlength{\LTleft}{0pt}\setlength{\LTright}{\fill}
\begin{longtable}{@{}p{0.2200\linewidth}p{0.1750\linewidth}p{0.1750\linewidth}p{0.1750\linewidth}p{0.1750\linewidth}@{}}
\caption{Negative odd roots at every nonunit inner length: parameter values.}\label{tab:parameters:q5-c080:4}\\
\toprule
Entry & eta & \(E_0/d\) upper & Extra norm cost divided by d & Final logarithmic upper divided by d at \(k=1\) \\
\midrule\endfirsthead
\toprule
Entry & eta & \(E_0/d\) upper & Extra norm cost divided by d & Final logarithmic upper divided by d at \(k=1\) \\
\midrule\endhead
\bottomrule\endfoot
Every \(L=2\); also \(b=2,\allowbreak  L=3\) & \(11/12\); & \(3/4\); & \(53/4\); & \(-28/3\). \\
\addlinespace[2pt]
\(b=2,\allowbreak L=4\) & \(17/16\); & \(1/4\); & 15; & \(-97/12\). \\
\addlinespace[2pt]
\(b=2,\allowbreak L\ge 5\) & \(17/16\); & 0; & 15; & \(-25/3\). \\
\addlinespace[2pt]
\(b\ge 3,\allowbreak L=3\) & \(11/10\); & \(-1/10\); & \(309/20\); & \(-479/60\). \\
\addlinespace[2pt]
\(b\ge 3,\allowbreak L=4\) & \(15/16\); & 0; & \(27/2\); & \(-59/6\). \\
\addlinespace[2pt]
\(b\ge 3,\allowbreak L\ge 5\) & \(6/7\); & 0; & \(351/28\); & \(-907/84\). \\
\end{longtable}
\endgroup


The scalar and norm calculations for the table are as follows.

For the first row, (Q80.3.5) has \(D_2/(B^2d)<2.\) Conversion gives a
one-row norm at most \(2\theta=13/16<(11/12)^2.\) The weakest original
scalar is -d/8; conversion costs less than 7d/8. Their sum is
at most 3d/4. The extra norm cost is
\(16(11/16+9/64)d=53d/4.\) The full sum is
\(3/4+53/4+2/3-24=-28/3.\)

For \(b=2,\allowbreak  L=4,\allowbreak \) (Q80.3.5) gives \(D_2/(B^2d)\le 217/81.\) The one-row norm
is at most \(\theta(217/81)=2821/2592<(17/16)^2.\) Conversion costs
at most 1519d/1296, while the original scalar is at most -d.
Thus its scalar upper is \(223d/1296<d/4.\) The extra norm cost is
\(16(51/64+9/64)d=15d.\) This gives \(1/4+15+2/3-24=-97/12.\)

For \(b=2,\allowbreak  L\ge 5,\allowbreak \) (Q80.3.6) is exactly at most (\(17/16)^2\) and its original
one-row scalar is nonpositive. No conversion cost is added.
The same extra norm cost 15d gives \(15+2/3-24=-25/3.\)

For \(b\ge 3,\allowbreak  L=3,\allowbreak \) (Q80.3.6) gives \(481/400<(11/10)^2,\allowbreak \) with the complete
original scalar at most -d/10. The extra norm cost is
\(16(33/40+9/64)d=309d/20.\) The resulting sum is
-\(1/10+309/20+2/3-24=-479/60.\)

For \(b\ge 3,\allowbreak  L=4,\allowbreak \) the squared norm is at most \(27/32<(15/16)^2\)
and the original scalar is nonpositive. The extra norm cost
\(16(45/64+9/64)d=27d/2\) gives \(27/2+2/3-24=-59/6.\)

For \(b\ge 3,\allowbreak  L\ge 5,\allowbreak \) the squared norm is at most \(18/25<(6/7)^2\)
and the original scalar is nonpositive. The extra norm cost
\(16(9/14+9/64)d=351d/28\) gives \(351/28+2/3-24=-907/84.\)

Every table entry is at most -\(479/60.\) In particular, the close
second row is -\(97/12=-485/60,\allowbreak \) with surplus \(1/10.\) The first and
third give losses \(560/60\) and \(500/60\); the fifth gives \(590/60.\)
For the last comparison, \(907/84-479/60=1182/420>0.\)
All \(k>1\) are stronger by 24(k-1)d. This proves the ordinary part
of Theorem T085.2 throughout \(0\le d\le 1/3.\)

The short-three exception \(u_1=3\) is excluded here: its first-pair loss is not \(3d\), so (Q80.1.3) and this table do
not establish that boundary. No conclusion is drawn from this
sufficient payment about whether that boundary is true or false.
